\begin{lemma}
\label{lemma-compute-tensor-perfect}
Assumptions and notation as in Lemma \ref{lemma-tensor-perfect}.
Then there are functorial isomorphisms
$$
H^i(S, K \otimes^\mathbf{L}_{\mathcal{O}_S} \mathcal{F})
\longrightarrow
H^i(X, E \otimes_{\mathcal{O}_X}^\mathbf{L}
(\mathcal{G}^\bullet \otimes_{\mathcal{O}_X} f^*\mathcal{F}))
$$
for $\mathcal{F}$ quasi-coherent on $S$
compatible with boundary maps (see proof).
\end{lemma}

\begin{proof}
We have
$$
\mathcal{G}^\bullet \otimes_{\mathcal{O}_X}^\mathbf{L} Lf^*\mathcal{F} =
\mathcal{G}^\bullet \otimes_{f^{-1}\mathcal{O}_S}^\mathbf{L} f^{-1}\mathcal{F} =
\mathcal{G}^\bullet \otimes_{f^{-1}\mathcal{O}_S} f^{-1}\mathcal{F} =
\mathcal{G}^\bullet \otimes_{\mathcal{O}_X} f^*\mathcal{F}
$$
the first equality by
Cohomology, Lemma \ref{cohomology-lemma-variant-derived-pullback},
the second as $\mathcal{G}^n$ is a flat $f^{-1}\mathcal{O}_S$-module, and
the third by definition of pullbacks. Hence we obtain
\begin{align*}
H^i(X, E \otimes^\mathbf{L}_{\mathcal{O}_X}
(\mathcal{G}^\bullet \otimes_{\mathcal{O}_X} f^*\mathcal{F}))
& =
H^i(X, E \otimes^\mathbf{L}_{\mathcal{O}_X} \mathcal{G}^\bullet
\otimes_{\mathcal{O}_X}^\mathbf{L} Lf^*\mathcal{F}) \\
& =
H^i(S,
Rf_*(E \otimes^\mathbf{L}_{\mathcal{O}_X} \mathcal{G}^\bullet
\otimes^\mathbf{L}_{\mathcal{O}_X} Lf^*\mathcal{F})) \\
& =
H^i(S, Rf_*(E \otimes^\mathbf{L}_{\mathcal{O}_X} \mathcal{G}^\bullet)
\otimes^\mathbf{L}_{\mathcal{O}_S} \mathcal{F}) \\
& =
H^i(S, K \otimes^\mathbf{L}_{\mathcal{O}_S} \mathcal{F}) 
\end{align*}
The first equality by the above, the second by Leray
(Cohomology, Lemma \ref{cohomology-lemma-before-Leray}), and
the third equality by Lemma \ref{lemma-cohomology-base-change}.
The statement on boundary maps means the following: Given a short
exact sequence $0 \to \mathcal{F}_1 \to \mathcal{F}_2 \to \mathcal{F}_3 \to 0$
of quasi-coherent $\mathcal{O}_S$-modules, the isomorphisms fit into
commutative diagrams
$$
\xymatrix{
H^i(S, K \otimes^\mathbf{L}_{\mathcal{O}_S} \mathcal{F}_3)
\ar[r] \ar[d]_\delta &
H^i(X, E \otimes^\mathbf{L}_{\mathcal{O}_X}
(\mathcal{G}^\bullet \otimes_{\mathcal{O}_X} f^*\mathcal{F}_3))
\ar[d]^\delta \\
H^{i + 1}(S, K \otimes^\mathbf{L}_{\mathcal{O}_S} \mathcal{F}_1)
\ar[r] &
H^{i + 1}(X, E \otimes^\mathbf{L}_{\mathcal{O}_X}
(\mathcal{G}^\bullet \otimes_{\mathcal{O}_X} f^*\mathcal{F}_1))
}
$$
where the boundary maps come from the distinguished triangle
$$
K \otimes^\mathbf{L}_{\mathcal{O}_S} \mathcal{F}_1 \to
K \otimes^\mathbf{L}_{\mathcal{O}_S} \mathcal{F}_2 \to
K \otimes^\mathbf{L}_{\mathcal{O}_S} \mathcal{F}_3 \to
K \otimes^\mathbf{L}_{\mathcal{O}_S} \mathcal{F}_1[1]
$$
and the distinguished triangle in $D(\mathcal{O}_X)$ associated to
the short exact sequence
$$
0 \to
\mathcal{G}^\bullet \otimes_{\mathcal{O}_X} f^*\mathcal{F}_1 \to
\mathcal{G}^\bullet \otimes_{\mathcal{O}_X} f^*\mathcal{F}_2 \to
\mathcal{G}^\bullet \otimes_{\mathcal{O}_X} f^*\mathcal{F}_3 \to 0
$$
of complexes of $\mathcal{O}_X$-modules.
This sequence is exact because $\mathcal{G}^n$ is flat over $S$.
We omit the verification of the commutativity of the displayed diagram.
\end{proof}
